% ============================================================
% INTRODUÇÃO GERAL
% ============================================================
\chapter*{Introdução Geral}
\addcontentsline{toc}{chapter}{Introdução Geral}

\section*{Por que este livro, por que agora}

O ambiente de negócios brasileiro em 2024–2025 é definido por \textbf{quatro choques simultâneos}:
\begin{enumerate}
    \item \textbf{Macroeconômico:} Selic 10,75\% (dez/2024), IPCA 4,6\% (12m), dívida/PIB 76\%, reforma tributária em transição (EC 132/2023). Custo de capital real \textasciitilde 4,5\% a.a. — filtragem implacável de projetos.
    \item \textbf{Tecnológico:} IA generativa (ChatGPT nov/2022 → adoção empresarial 34\% — CGI.br 2024), cloud 68\%, automação/hiperautomação, gêmeos digitais. \textit{AI-first} deixa de ser opção para ser condição de competitividade.
    \item \textbf{Regulatório/Social:} LGPD (ANPD ativa), Lei Anticorrupção (Decreto 11.129/2022), Marco Legal IA (PLC 2338/2023 tramitando), ESG mandatório (CVM Res. 193, ISSB S1/S2, CSRD/ESRS se Europa), diversidade obrigatória (Lei 14.611/2023, B3 Mulheres na Liderança).
    \item \textbf{Climático:} SEEG 2024 — 2,3 GtCO₂e (agro 28\%, energia 25\%, LULUCF 22\%), eventos extremos (RS 2024), Net Zero 2050 (SBTi), taxonomia verde BCB, mercado carbono (PL 182/2024).
\end{enumerate}

Empresas brasileiras que navegam esta tempestade compartilham \textbf{cinco atributos}:
\begin{enumerate}
    \item \textbf{Clareza estratégica:} sabem \textit{onde jogar} e \textit{como vencer} (Porter), com \textit{capacidades dinâmicas} (Teece) para reconfigurar-se.
    \item \textbf{Disciplina de capital:} ROIC > WACC consistentemente, alocação por \textit{três horizontes} (McKinsey), \textit{rolling forecast} substituindo orçamento estático.
    \item \textbf{Excelência operacional:} Lean/Seis Sigma no DNA, resiliência cadeia por design (visibilidade Tier 3, dual sourcing), digitalização com ROI (manutenção preditiva, roteirização IA).
    \item \textbf{Gestão de talentos como ativo:} pipeline liderança (Charan/Drotter/Noel), 9-Box, OKR cascade, cultura experimentação, DEIP genuíno.
    \item \textbf{Governança \& confiança:} conselho independente, compliance efetivo (ISO 37301/37001, Three Lines), autoregulação, relato transparente (GRI/ISSB/TCFD).
\end{itemize}

\section*{Estrutura e como usar}

\begin{itemize}
    \item \textbf{Executivos/C-level:} Leiam \textit{Conclusões} de cada capítulo (imperativos gerenciais) + \textit{Figuras/Tabelas} (decisões visuais) + \textit{Introdução Geral} + \textit{Considerações Finais}.
    \item \textbf{Gestores/Coordenadores:} Foquem nos \textit{Eixos de Fundamentação} (frameworks aplicáveis) + \textit{Metodologia} (como implementar) + \textit{Artefatos} (templates adaptáveis).
    \item \textbf{Estudantes/Pesquisadores:} \textit{Fundamentação Teórica} (referências seminal + 2020–2025) + \textit{Agenda de Pesquisa} (gaps identificados) + \textit{Referências ABNT} (ponto de partida bibliográfico).
    \item \textbf{Consultores/Auditores:} \textit{Matrizes de Diagnóstico} (Compliance 10 Pilares, Kraljic, Kraljic Adaptada, Kraljic, Salience, Bow-Tie, Heat Map, RAS, AI Maturity, GenAI Portfolio, Responsible AI) + \textit{Checklists} implícitos nas figuras.
\end{itemize}

\section*{Metodologia transversal}

\begin{enumerate}
    \item \textbf{Levantamento bibliográfico global (Etapa 2):} 14 capítulos $\times$ 4–6 eixos $\times$ 20+ referências cada = 1.200+ referências mapeadas (seminais + 2020–2025). Fontes: Scopus, Web of Science, OpenAlex, Semantic Scholar, bases institucionais (McKinsey, BCG, Gartner, Deloitte, PwC, KPMG, WEF, OECD, IBGC, FGV, CNI, IBGE, BCB, CVM, B3, CGI.br, ANPD, ISO, COSO, IIA, PRI, GRI, ISSB, SBTi, TCFD, NGFS).
    \item \textbf{Cases integradores (3 por capítulo):} Magazine Luiza (ecossistema digital, varejo+fintech+logística+marketplace), Weg (industrial global, manufatura+engrenharia+energia+mobilidade), Nubank (nascida digital, banking+invest+insurance+credit, cultura owner/OKR/MLOps) + 1 case setorial por capítulo (Embraer, Natura, Ambev, Petrobras, JBS, Hospital Einstein, Banco do Brasil, Vale).
    \item \textbf{Dados reais 2024–2025:} Todos quantitativos, fontes primárias citadas, benchmarks setoriais.
    \item \textbf{Artefatos visuais originais (56):} Desenhados em TikZ/pgfplots (LaTeX nativo), vetoriais, editáveis, consistentes na identidade visual (cores ecnavy/ecgold/ecteal/eccoral).
    \item \textbf{Validação:} 2 professores estratégia (FGV/EAESP), 1 C-level varejo, 1 consultor BCG/McKinsey — \textit{pendente no momento da redação}.
\end{enumerate}

\section*{Pressupostos e limitações}

\begin{itemize}
    \item Foco em \textbf{grandes e médias empresas brasileiras} (listadas ou pré-IPO). PMEs e terceiro setor merecem obra dedicada.
    \item Recorte temporal: \textbf{dados até dez/2024}, projeções 2025 baseadas em guidance oficial e consenso de mercado (Focus BCB, FMI, OCDE).
    \item Referências 2020–2025 priorizam \textbf{relatórios institucionais públicos} e \textbf{artigos de alto impacto} — DOIs/páginas omitidos para fluidez; validar na redação final.
    \item Cases baseados em \textbf{informações públicas} (relatórios anuais, earnings calls, entrevistas publicadas, tech blogs). Dados internos não divulgados não foram usados.
    \item \textbf{IA generativa} usada como ferramenta de apoio (levantamento, estruturação, redação, código TikZ, revisão) — autor manteve controle intelectual integral (conforme Declaração de Uso de IA).
\end{itemize}

\section*{Convite ao leitor}

A excelência corporativa não se decreta — \textit{se constrói}. Este livro oferece mapas, bússolas e ferramentas. A jornada é de cada organização. Que as próximas páginas provoquem perguntas melhores, decisões mais corajosas e resultados mais duradouros.

\vspace{0.5cm}
\noindent\textit{Boa leitura e boa gestão.}
\end{flushright}